\documentclass{beamer}
\usetheme{Boadilla}
\usepackage{xcolor}
\usepackage{amsmath}
\usepackage{amssymb}
\definecolor{kkcolor}{RGB}{61, 66, 107}
\usecolortheme[named=kkcolor]{structure}
\usepackage[
backend=biber,
style=apa,
]{biblatex}
\newcommand{\figpath}{../../3. Output/figures/}
\newcommand{\tabpath}{../../3. Output/tables/}

\title[The Decision to Hold Back]{The Decision to Hold Back: Experimental Evidence on Teachers’ Repetition Biases}
\author[Garijo Campos, Ignacio]{Ignacio Garijo Campos \\
		Antonio Alfonso \\
		Pablo Brañas \\
		Diego Jorrat \\
		Ángel Sánchez 
}
\institute[]{Universidad Loyola Andalucía}

\AtBeginSection[]
{
	\begin{frame} 

		\tableofcontents[currentsection]
	\end{frame}
}

\begin{document}
	

		\maketitle

	
	
	\section{Introduction} \label{intro}
	
	\begin{frame}{\nameref{intro}: Grade Repetition}
		\begin{itemize}
			\item Grade repetition can be seen as a \textbf{double punishment}: students face academic failure and are separated from their peers.
			\item This separation may negatively affect \textbf{self-esteem} and perceptions of ability.
			\item In psychology, this relates to \textbf{learned helplessness}: students stop trying because they believe success is impossible regardless of effort.
			\item Overall, the literature finds \textbf{little or no positive effects} of grade repetition on test scores, skills or dropout rates
		\end{itemize}
	\end{frame}
	
	
	\section{Our approach}\label{exp}
	
	\begin{frame}{\nameref{exp}}
		\begin{itemize}
			\item Our novel approach: we focus on teachers and their `harshness' level
			\item We carry out a survey experiment in Spain, where the grade repetition rate more than doubles that of the OECD (9\% and 22\%)
			\item Double objective: 
				\begin{enumerate}
					\item Characterizing `harsh' teachers 
					\item Studying how teachers behave under different policy circumstances
				\end{enumerate}
		\end{itemize}
	\end{frame}
	
	\begin{frame}{\nameref{exp}: What is harshness?}
		
		Socioeconomic questionnaire $\rightarrow$ Experimental section $\rightarrow$ Behavioral questionnaire		
		 \begin{table}[!ht]
			\centering 
			\begin{tabular}{|p{7cm}|*{3}{c|}}
				\hline
				& Option 1 & Option 2 \\
				\hline
				Gender & Male & Female \\
				\hline
				Disadvantaged background & Yes & No  \\
				\hline
				Failed 3 or more subjects & Yes & No  \\
				\hline
				Lacks mathematical and/or linguistic abilities & Yes & No  \\
				\hline
				Has committed a serious or very serious infraction or has been suspended & Yes & No\\
				\hline
				Frequently misses class without justification & Yes & No \\
				\hline
			\end{tabular}
			\label{tab:cards}
		\end{table}
	\end{frame}
	
	\begin{frame}{\nameref{exp}: Measures of harshness}
		We estimate 4 different measures, but due to their similarity, only show results for one: 
		
		\begin{enumerate}
			\item Number of repetitions
			\item \textbf{Average deviation with respect to peers} ($\simeq$ 320)
			\item Average deviation with respect to predicted probability of repetition
			\item Average deviation with respect to predicted probability of repetition in a more complete model
		\end{enumerate}
		
		\begin{figure} 
			\centering
			\caption{Average deviation with respect to peers}
			\includegraphics[width=0.6\textwidth]{\figpath histogram_hb.jpeg}
		\end{figure}
	\end{frame}
	
	\section{Characterizing harsh teachers}\label{charac}
	
	\begin{frame}{\nameref{charac}: the model}
		\begin{figure} 
			\centering
			\caption{Coefficients of linear regression on harshness}
			\includegraphics[width=\textwidth]{\figpath coefs.jpeg}
		\end{figure}
	\end{frame}

\end{document}